% ch6_a 更一般的黑洞 (书页 238-250 / p253-p265)
% 第6章 More General Black Holes
\chapter{更一般的黑洞}

\section{黑洞动物园}

% ---- 书页 238 / p253 ----
Birkhoff 定理保证了 Schwarzschild 度规是广义相对论中唯一的球对称真空解。这并不太令人惊讶，因为它让我们想起电磁学中的情形：在无电荷的区域里，唯一的球对称场位形就是 Coulomb 场。走出球对称性之后，可能的引力场就五花八门、不胜枚举。比如对一颗像地球这样的行星，其外部场将依赖于表面上所有各种山脉与山谷的密度和轮廓。我们可以设想把度规按多极矩（multipole moments）展开；要精确描述这个场，就得给定无穷多个系数。

因此，黑洞并不具备这种性质，这一点或许反倒会让人感到几分意外。稳态黑洞解只有寥寥几个，每个解由少数几个参量描述。参量的具体集合取决于我们的理论中包含哪些物质场；如果电磁学是唯一的长程非引力场，我们就有一条\textbf{无毛定理}（no-hair theorem）：

与电磁学相耦合的广义相对论的稳态、渐近平直黑洞解，只要在事件视界之外无奇点，就完全由质量、电荷与磁荷以及角动量这几个参量刻画。

稳态解之所以特别令人感兴趣，是因为我们预期它们是引力坍缩的终态。另一种可能是某种振荡的位形，但随着能量通过引力辐射的发射而损耗，振荡终将被阻尼；实际上，典型的演化会相当迅速地趋于稳态位形。

我们说的是“一条”无毛定理，而不是“那条”无毛定理，因为这一结果不仅依赖于广义相对论，还依赖于我们理论中的物质内容。在粒子物理的标准模型中，电磁学是唯一的长程场，上面的定理适用；但对别的种类的场来说，可能存在其他种类的毛发。\footnote{相关讨论见 M. Heusler, “Stationary Black Holes: Uniqueness and Beyond,” \emph{Living Rev.\ Relativity} \textbf{1} (1998): 6, http://www.livingreviews.org/Articles/Volume1/1998-6heusler/.} 甚至还有人找到了静态（无转动）的、轴对称却不完全球对称的黑洞解的例子。

% ---- 书页 239 / p254 ----
然而核心结论没有改变：黑洞解由极少数参量刻画，而不像比如说一颗行星那样，需要可能无穷多的参量来刻画。

正如我们将在本章末尾以及在第 9 章再次讨论的那样，无毛性质导致了一个令人困惑的局面。在大多数物理理论中，我们希望有一个良定义的初值问题，从而关于任一时刻之态的信息都可以用来预言（或回溯）任何其他时刻的态。这样一来，由运动方程的一个解联系起来的任何两个态，就应当需要同样多的信息来给定。但在 GR 中，事情似乎是这样的：我们可以取一团非常复杂的物质，让它坍缩成黑洞，最终得到的位形却完全由质量、电荷和自旋来描述。在经典 GR 中这或许还不会太困扰我们，因为可以认为这些信息只是藏在事件视界背后，而非真正丢失。但一旦把量子场论考虑进来，我们就会发现黑洞会蒸发并最终消失，信息看起来就真的丢失了。可以想见，导致蒸发的出射霍金辐射或许以某种方式编码了关于原来是用哪个态来造出黑洞的信息，但这一切如何发生完全不清楚。在许多人看来，理解这个“信息丢失佯谬”（information loss paradox）是构建一门合理的量子引力理论的关键一步。

不过在本章中，我们将只限于讨论经典 GR。我们先对黑洞的性质作一些一般性的讨论，尤其是事件视界和 Killing 视界的性质。这个话题可以相当微妙和技术化，我们在这里的方针是力图传达主要思想，而不在定义和定理证明上讲求严格。然后我们讨论带电（Reissner–Nordström）黑洞和转动（Kerr）黑洞所对应的具体解；与我们的一贯做法一致，我们不会仔细梳理构造最大延拓时空所需的那些坐标重新定义，而是径直画出相应的共形图。想进一步了解细节的读者，可以查阅 Townsend 的综述文章\footnote{P.K. Townsend, “Black Holes: Lecture Notes,” http://arxiv.org/abs/gr-qc/9707012.}，以及 Hawking 与 Ellis (1973) 和 Wald (1984) 的书；本章大量取材于这些文献。

\section{事件视界}

黑洞的特点在于：你可以进入它，却永远无法出来。因此，黑洞最重要的特征其实并不是位于中心的奇点，而是位于边界的事件视界。事件视界是一个超曲面，它把能用类时路径与无穷远相连通的时空点，同那些不能连通的时空点分隔开来。要在实践中理解这句话的含义，我们需要更仔细地想一想“无穷远”指的是什么。在广义相对论中，时空的整体结构可以呈现出多种不同的形态，相应地也就有各不相同的无穷远概念。但要思考真实宇宙中的黑洞，我们
% ---- 书页 240 / p255 ----
其实并不真正关心无穷远处发生的事情；我们把无穷远当作“远在黑洞之外”的代名词，并且设想离洞足够远的时空可以用 Minkowski 空间来近似。

%FIG{6.1}: {渐近平直时空是指：其共形图中的无限远与 Minkowski 时空的无限远一致，包括未来类光无限远 $\mathscr{I}^{+}$、类空无限远 $i^{0}$ 和过去类光无限远 $\mathscr{I}^{-}$。未来事件视界是 $\mathscr{I}^{+}$ 的过去的边界。虚线区域代表时空的其余部分，在不同的例子中它可以呈现多种不同的形式。}
正如第 5 章提到的，在无穷远处看起来像 Minkowski 时空的时空被称为渐近平直（asymptotically flat）时空。这个概念的含义在一幅诸如图 6.1 那样的共形图中一目了然。从附录 H 中关于 Minkowski 时空共形图的讨论我们知道，共形无限远分为五块：未来与过去类时无限远 $i^{\pm}$、未来与过去类光无限远 $\mathscr{I}^{\pm}$，以及类空无限远 $i^{0}$。如果一个时空（或时空的一个区域）的 $\mathscr{I}^{\pm}$ 和 $i^{0}$ 具有与 Minkowski 时空相同的结构，它就是渐近平直的；类时无限远则不是必需的。这样的时空一般具有图 6.1 所示的形态。

有了这幅图像，我们应该如何理解未来事件视界就清楚了：它是类时曲线无法再逃逸到无穷远的那一侧边界曲面。回忆一个区域的因果过去 $J^{-}$，是指从该区域出发、沿指向过去的类时路径所能到达的所有点的集合；事件视界于是可以等价地定义为 $J^{-}(\mathscr{I}^{+})$——即未来类光无限远的因果过去——的边界。（事件视界其实应当是这个集合的\emph{闭包}（closure）的边界，不过我们在这里不追求严格。）对过去视界也有类似的定义。正如在对 Schwarzschild 作最大延拓的情形中所看到的，一个时空中可以不止有一个渐近平直区域，相应地也就不止有一个事件视界。

从定义可以看出，事件视界是一个类光超曲面。类光超曲面的性质在附录 D 中讨论过；这里我们回顾一下其主要特点。一个超曲面 $\Sigma$ 可以由某个函数 $f(x)$ 通过 $f(x)=$ 常数来定义。梯度 $\partial_\mu f$ 与 $\Sigma$ 正交；若法矢量是类光的，就称这个超曲面为类光超曲面，此时法矢量同时也切于 $\Sigma$。类光超曲面可以看成是一族类光测地线 $x^{\mu}(\lambda)$ 的集合，这些测地线称为该超曲面的生成元（generators）。这些测地线的切矢量 $\xi^{\mu}$ 与法矢量成比例，

% ---- 书页 241 / p256 ----
\begin{equation*}
\xi^\mu = \frac{dx^\mu}{d\lambda} = h(x)g^{\mu\nu}\partial_\nu f,
\tag{6.1}
\end{equation*}
因而也可以充当超曲面的法矢量。我们可以选取函数 $h(x)$，使得这些测地线是仿射参量化的，于是切矢量将满足
\begin{equation*}
\xi_\mu\xi^\mu = 0, \qquad \xi^\mu\nabla_\mu\xi^\nu = 0.
\tag{6.2}
\end{equation*}
对于未来事件视界，其生成元可以在过去终结（例如，当黑洞由恒星坍缩形成时），但总会一直延伸到无穷远的未来（未来与过去互换后亦然）。

由于事件视界是一个整体性的概念，如果有人递给你一个用任意坐标写出的度规，要真正找到它的视界位置可能相当困难。幸运的是，本章我们要打交道的都是相当特殊的度规——稳态、渐近平直，而且含有拓扑为球面的事件视界。在这类时空中，存在一些方便的坐标系，使得辨认事件视界有一种简单办法。对 Schwarzschild 解来说，事件视界是光锥“倾倒”的地方，使得 $r=2GM$ 成为一个类光曲面，而不像在大 $r$ 处那样 $r=$ 常数是类时曲面。光锥倾倒显然是一个依赖坐标的概念（比如在 Kruskal 坐标中就不会发生），但我们关心的那些度规都容许类似的构造。稳态度规有一个渐近类时的 Killing 矢量 $\partial_t$，我们可以调整度规分量使之与时间无关（$\partial_t g_{\mu\nu}=0$）。在 $t=$ 常数的超曲面上，我们可以选取坐标 $(r,\theta,\phi)$，使得无穷远处的度规看起来就像球坐标下的 Minkowski 空间。$r=$ 常数的超曲面在 $r\to\infty$ 处将是拓扑为 $S^{2}\times\mathbf{R}$ 的类时柱面。现在设想我们聪明地选好了坐标，使得当 $r$ 从无穷远逐渐减小时，$r=$ 常数的超曲面始终保持类时，直到某个固定的 $r=r_{\mathrm{H}}$ 处，曲面处处类光。（在不那么聪明的坐标中，$r=$ 常数的超曲面会在某些 $\theta$ 和 $\phi$ 值处变成类光或类空，而在其他值处仍是类时。）这显然就代表一个事件视界，因为跨入 $r<r_{\mathrm{H}}$ 区域的类时路径永远无法逃回无穷远。确定 $r=$ 常数超曲面在何处变为类光很容易：$\partial_\mu r$ 是这种超曲面的一个法 1-形式，其模长为
\begin{equation*}
g^{\mu\nu}(\partial_\mu r)(\partial_\nu r) = g^{rr}.
\tag{6.3}
\end{equation*}
我们要找的正是这个 1-形式的模长为零的地方；于是，在我们所描述的坐标中，事件视界 $r=r_{\mathrm{H}}$ 就是 $g^{rr}$ 由正变负的那个超曲面，
\begin{equation*}
g^{rr}(r_{\mathrm{H}}) = 0.
\tag{6.4}
\end{equation*}

% ---- 书页 242 / p257 ----
这条判据对 Schwarzschild 解显然有效，因为它的 $g^{rr}=1-2GM/r$。后面我们给出 Reissner–Nordström 解和 Kerr 解时，也将采用类似地适配于视界的坐标。

我们之所以把事件视界看得如此重要，是因为在广义相对论中它们几乎是不可避免的。这一结论由两个有趣的结果串联而来：奇点几乎不可避免，而且奇点藏在事件视界背后。当然，两个结果都在各自的一组适当假设下才成立；构造没有奇点的时空并不难（Minkowski 空间就是一例），找到没有视界的奇点也不算难（下面讨论带电黑洞时就会看到）。但我们相信，“通有”（generic）的解都会有藏在视界背后的奇点。

奇点的普遍存在由 Hawking 和 Penrose 的\textbf{奇点定理}（singularity theorems）保证。在这些定理被证明之前，人们还可以指望，坍缩到 Schwarzschild 奇点只是球对称性的一个人为假象，典型的几何将保持无奇点（例如在 Newton 引力中就是这样）。但 Hawking–Penrose 定理表明，坍缩一旦达到某个程度，向奇点的演化就不可避免。我们知道存在奇点的途径是测地不完备（geodesic incompleteness）——存在某条测地线，它无法在流形内延伸，却在仿射参量的有限值处终结。而我们知道坍缩已经达到有去无回之点的标志，是一个\textbf{陷俘面}（trapped surface）的出现。要理解什么是陷俘面，先想象 Minkowski 空间中的一个二维球面，它嵌在一个等时切片里，取为到原点径向距离固定的一些点。如果我们追踪从这个空间球面射向时空的类光线，一族（指向内部）将描绘出一组不断收缩的球面，而另一族（指向外部）将描绘出一组不断胀大的球面。但对 Schwarzschild 几何中半径固定为 $r<2GM$ 的球面来说，情况就不是这样了；在事件视界之内，从这个球面发出的两组类光线都会演化到更小的 $r$ 值（因为 $r$ 是类时坐标），从而演化到更小的面积 $4\pi r^{2}$。这就是陷俘面的含义：一个紧致、类空的二维子流形，其性质是在该子流形上处处，指向未来的出射光线沿两个方向都\emph{汇聚}。（“汇聚”的正式定义是：膨胀 $\theta$——正如附录 F 中关于测地线汇的讨论所描述的——为负。）

有了这些定义，我们就可以给出一个奇点定理的例子。

令 $M$ 是一个流形，带有通有的度规 $g_{\mu\nu}$，满足 Einstein 方程并施加强能量条件。如果 $M$ 中存在一个陷俘面，则必定要么存在一条闭合类时曲线，要么存在一个奇点（表现为某条不完备的类时或类光测地线）。

% ---- 书页 243 / p258 ----
这里所谓“通有度规”，指的是对类时和类光测地线都满足\textbf{普适条件}（generic condition）。对类时测地线，普适条件是说：每条以 $U^{\mu}$ 为切矢量的测地线，都必须至少有一个点使得 $R_{\alpha\beta\gamma\delta}U^{\alpha}U^{\delta}\neq 0$；对类光测地线，普适条件是说：每条以 $k^{\mu}$ 为切矢量的测地线，都必须至少有一个点使得 $k_{[\alpha}R_{\beta]\gamma\delta[\epsilon}k_{\xi]}k^{\gamma}k^{\delta}\neq 0$。这些花哨的条件不过是用来排除那些曲率在某些方向上一致为零的特殊度规。

奇点定理有许多种形式，各自从不同的一组假设出发。这个故事的寓意似乎是：广义相对论中典型的含时解通常终于奇点。（或者始于奇点；某些定理暗示了宇宙学奇点的存在，比如大爆炸。）这对 GR 来说算是个问题，因为这个理论其实并不真正适用于奇点本身，奇点的存在因此代表着描述的不完备性。对待这个问题的传统态度，是寄希望于众人苦苦追寻的量子引力理论，能以某种方式化解经典 GR 的奇点。

与此同时，我们还可以从“奇点藏在事件视界背后”这一想法中获得安慰。这一信念概括在\textbf{宇宙监督假设}（cosmic censorship conjecture）之中：

在满足主能量条件的渐近平直时空中，从通有的、初始非奇异的态出发的引力坍缩，不能形成裸奇点。

\textbf{裸奇点}（naked singularity）是指信号可以从那里到达 $\mathscr{I}^{+}$ 的奇点；也就是说，是没有藏在事件视界背后的奇点。注意，这个假设讲的是裸奇点的\emph{形成}，而不是它们的\emph{存在}；当然有一些解，其类空裸奇点存在于过去（比如 Schwarzschild 白洞），或者类时奇点存在于所有时间（比如下面要讨论的超极端带电黑洞）。宇宙监督假设至今未被证明，尽管人们已经花了大量力气去寻找有说服力的反例，却一无所获。要求初始数据在某种意义上“通有”这一点很重要，因为数值实验表明，经过精细调节的初始条件能够产生裸奇点。对某种形式的宇宙监督假设给出严格证明，仍然是经典广义相对论悬而未决的突出问题之一。\footnote{关于宇宙监督假设的评述见 R.M. Wald, “Gravitational Collapse and Cosmic Censorship,” http://arxiv.org/abs/gr-qc/9710068.}

宇宙监督（或某些等价假设）的一个推论是：经典黑洞从不缩小，只会越长越大。黑洞的大小由事件视界的面积来量度，这里指的是事件视界与一个类空切片之交的空间面积。于是我们有了 Hawking 的\textbf{面积定理}（area theorem）：

假设弱能量条件与宇宙监督成立，则渐近平直时空中未来事件视界的面积不减。

% ---- 书页 244 / p259 ----
对 Schwarzschild 黑洞，面积单调地依赖于质量，因此这个定理意味着 Schwarzschild 黑洞的质量只能增加。但对转动的黑洞就不再如此了；面积依赖于质量和角动量的某种组合，而且下面我们会看到，通过减小黑洞的自旋，我们真的可以从中提取能量。我们也可以通过量子力学的霍金辐射来减小黑洞的质量；这可以追溯到弯曲时空中的量子场论可以违反弱能量条件这一事实。

\section{Killing 视界}

在 Schwarzschild 度规中，Killing 矢量 $K=\partial_t$ 在事件视界处由类时变为类空。一般地，如果某个 Killing 矢量场 $\chi^{\mu}$ 沿某个类光超曲面 $\Sigma$ 处处类光，我们就说 $\Sigma$ 是 $\chi^{\mu}$ 的一个\textbf{Killing 视界}（Killing horizon）。注意矢量场 $\chi^{\mu}$ 必然与 $\Sigma$ 正交，因为一个类光曲面不可能有两个线性无关的类光切矢量。

Killing 视界的概念在逻辑上独立于事件视界的概念，但在具有时间平移对称性的时空中，两者紧密相关。在某些合理的条件下（下面会明确给出），我们有如下分类：

稳态、渐近平直时空中的每一个事件视界 $\Sigma$，都是某个 Killing 矢量场 $\chi^{\mu}$ 的 Killing 视界。

如果时空是静态的，$\chi^{\mu}$ 就是代表无穷远处时间平移的 Killing 矢量场 $K^{\mu}=(\partial_t)^{\mu}$。

如果时空是稳态但非静态的，它将是轴对称的，并带有转动 Killing 矢量场 $R^{\mu}=(\partial_\phi)^{\mu}$，而 $\chi^{\mu}$ 将是线性组合 $K^{\mu}+\Omega_{\mathrm{H}}R^{\mu}$，其中 $\Omega_{\mathrm{H}}$ 为某个常数。

例如，下面我们将考察转动黑洞的 Kerr 度规，其中事件视界是转动与时间平移的 Killing 矢量的某个线性组合的 Killing 视界。在 Kerr 情形，$\partial_t$ 变为类光的那个超曲面实际上是类时的，因此并不是 Killing 视界。

让我们把这套分类方案究竟在什么条件下成立说得更精确些。\footnote{相关讨论见 R. M. Wald, “The thermodynamics of black holes,” \emph{Living Rev.\ Rel.} \textbf{4}, 6 (2001), http://arxiv.org/gr-qc/9912119.}Carter 已经证明，对静态黑洞，事件视界是 $K^{\mu}$ 的 Killing 视界；这是一个纯几何的事实，即使不引用 Einstein 方程也成立。在稳态情形，如果我们假设存在一个转动 Killing 场 $R^{\mu}$，它具有这样的性质：由 $K^{\mu}$ 和 $R^{\mu}$ 张成的二维平面与一族二维曲面正交，那么事件视界将是这两个 Killing 场的某个线性组合的 Killing 视界——这同样是纯几何的论证。反过来，如果我们只假设黑洞是稳态的，那么一般说来我们无法证明事件视界
% ---- 书页 245 / p260 ----
是轴对称的。给定 Einstein 方程加上对物质场的某些条件，Hawking 得以证明：任何稳态黑洞的事件视界必定是某个矢量场的 Killing 视界，而且这类视界必定要么稳态、要么轴对称。在本章剩下的部分里，我们将当作上述分类成立那样来叙述；然而，对物质场作假设是出了名的棘手，我们应当记住，原则上有可能找到既非静态也非轴对称的黑洞，对它们来说事件视界可能就不是 Killing 视界。

需要着重指出的一点是：虽然稳态渐近平直时空的事件视界通常是 Killing 视界，但也很容易找到与事件视界毫不相干的 Killing 视界。考虑用惯性坐标表示的 Minkowski 空间，$\mathrm{d}s^{2}=-\mathrm{d}t^{2}+\mathrm{d}x^{2}+\mathrm{d}y^{2}+\mathrm{d}z^{2}$；显然这个时空中没有任何事件视界。生成 $x$ 方向推动（boost）的 Killing 矢量为
\begin{equation*}
\chi = x\partial_t + t\partial_x,
\tag{6.5}
\end{equation*}
其模长为
\begin{equation*}
\chi_\mu\chi^\mu = -x^{2} + t^{2}.
\tag{6.6}
\end{equation*}
它在类光曲面
\begin{equation*}
x = \pm t,
\tag{6.7}
\end{equation*}
上变为类光，因此这些曲面就是 Killing 视界。把推动 Killing 矢量与平移和转动 Killing 矢量组合起来，我们就能让这些视界在流形中到处移动；Killing 视界俯拾皆是。当然，在更有意思的时空中，Killing 矢量场要少一些，相应的视界（如果有的话）也会有更大的物理意义。

对每一个 Killing 视界，我们都可以关联一个叫做\textbf{表面引力}（surface gravity）的量。考虑一个带有 Killing 视界 $\Sigma$ 的 Killing 矢量 $\chi^{\mu}$。因为 $\chi^{\mu}$ 是 $\Sigma$ 的法矢量，沿着 Killing 视界它满足测地线方程，
\begin{equation*}
\chi^\mu\nabla_\mu\chi^\nu = -\kappa\chi^\nu,
\tag{6.8}
\end{equation*}
等号右边的出现是因为 $\chi^{\mu}$ 的积分曲线未必是仿射参量化的。参量 $\kappa$ 就是表面引力；它在视界上为常数，除非在 Killing 矢量为零、$\kappa$ 可以变号的“分岔二维球面”（bifurcation two-sphere）处。（例如，在 Schwarzschild 解的 Kruskal 图中心就会发生这种情况。）利用 Killing 方程 $\nabla_{(\mu}\chi_{\nu)}=0$ 以及 $\chi_{[\mu}\nabla_\nu\chi_{\sigma]}=0$ 这一事实（因为 $\chi^{\mu}$ 与 $\Sigma$ 正交），可以直截了当地导出一个漂亮的表面引力公式：
\begin{equation*}
\kappa^{2} = -\frac{1}{2}(\nabla_\mu\chi_\nu)(\nabla^\mu\chi^\nu).
\tag{6.9}
\end{equation*}
右边的表达式要在视界 $\Sigma$ 上取值。鼓励你自己动手验证一下这个公式。

% ---- 书页 246 / p261 ----
与 Killing 视界相联系的表面引力原则上带有任意性，因为我们总可以用一个实常数去缩放某个 Killing 场而得到另一个 Killing 场。在静态、渐近平直的时空中，时间平移 Killing 矢量 $K=\partial_t$ 可以通过令
\begin{equation*}
K_\mu K^\mu(r\to\infty) = -1
\tag{6.10}
\end{equation*}
来归一化。这又随之确定了任何相关联的 Killing 视界的表面引力。如果我们处在稳态时空中，那里的 Killing 视界与时间平移和转动的某个线性组合相联系，那么固定 $K=\partial_t$ 的归一化也就固定了这个线性组合，因此表面引力仍然是唯一的。

$\kappa$ 为什么叫做“表面引力”，只有当时空是静态的时候才变得清楚。在这种情形下我们有如下解释：

在静态、渐近平直的时空中，表面引力就是视界附近一位静态观者的加速度，由无穷远处的一位静态观者测得。

为了弄清这句话的含义，我们先来考虑静态观者。所谓静态观者，指的是其四速度 $U^{\mu}$ 与时间平移 Killing 场 $K^{\mu}$ 成正比的观者，
\begin{equation*}
K^\mu = V(x)U^\mu.
\tag{6.11}
\end{equation*}
由于四速度归一化为 $U_\mu U^{\mu}=-1$，函数 $V$ 不过就是 Killing 场的模长，
\begin{equation*}
V = \sqrt{-K_\mu K^\mu},
\tag{6.12}
\end{equation*}
因而从 Killing 视界处的零一直取到无穷远处的 1。$V$ 有时被称为“红移因子”（redshift factor），因为它把静态观者所测得的一个光子的发射频率与观测频率联系起来。回忆一下，四动量为 $p^{\mu}$ 的光子的守恒能量是 $E=-p_\mu K^{\mu}$，而由四速度为 $U^{\mu}$ 的观者测得的频率将是 $\omega=-p_\mu U^{\mu}$。因此
\begin{equation*}
\omega = \frac{E}{V},
\tag{6.13}
\end{equation*}
而静态观者 1 发射的一个光子将被静态观者 2 观测到具有波长 $\lambda=2\pi/\omega$，
\begin{equation*}
\lambda_2 = \frac{V_2}{V_1}\lambda_1.
\tag{6.14}
\end{equation*}
特别地，在 $V=1$ 的无穷远处，我们将观测到波长 $\lambda_{\infty}=\lambda_{1}/V_{1}$。

现在我们转向“从无穷远处看到的加速度”这个概念。静态观者一般不会沿测地线运动；比如，粒子往往要落进黑洞，而不是悬停在黑洞旁边、保持空间坐标固定。

% ---- 书页 247 / p262 ----
我们可以把四加速度 $a^{\mu}=U^{\sigma}\nabla_{\sigma}U^{\mu}$ 用红移因子表示为
\begin{equation*}
a_\mu = \nabla_\mu\ln V,
\tag{6.15}
\end{equation*}
这一点你很容易验证。加速度的大小
\begin{equation*}
a = \sqrt{a_\mu a^\mu} = V^{-1}\sqrt{\nabla_\mu V\nabla^\mu V},
\tag{6.16}
\end{equation*}
在 Killing 视界处将趋于无穷——要让一个物体保持在静态轨迹上，需要无穷大的加速度。但无穷远处的观者探测到的加速度将被因子 $V$“红移”；这个结果正是表面引力。于是我们断言，
\begin{equation*}
\kappa = Va = \sqrt{\nabla_\mu V\nabla^\mu V},
\tag{6.17}
\end{equation*}
在视界 $\Sigma$ 上取值。你可以验证这个表达式与式 (6.9) 一致。表面引力是零（$V$）与无穷（$a$）的乘积，但通常取有限值。当我们说观测到的加速度被红移了，我们脑海中的图像是：把一根检验细绳从一个位于视界处的静态物体一直拉伸到无穷远处的观者，然后测量细绳在无穷远那一端的加速度。（值得花点时间看看，你能否把这个挥手带过的论证提升为某种更严格的东西。）

如果时空是稳态但非静态的，上面的考虑会在哪里出问题？我们仍然有一个渐近的时间平移 Killing 矢量 $K=\partial_t$，而且可以像式 (6.11) 那样，把四速度与 $K^{\mu}$ 平行的观者定义为稳态观者；红移仍由式 (6.14) 给出。问题在于，$K^{\mu}$ 不会在 Killing 视界处变为类光，而一般是在视界之外某个类时曲面上变为类光。这个 $K^{\mu}K_\mu=0$ 的地方被称为\textbf{静止极限面}（stationary limit surface）（有时也称“能层面”（ergosurface）），因为在这个曲面之内 $K^{\mu}$ 是类空的，因而任何观者都无法保持静止，哪怕这里还在事件视界之外。这样的观者不得不相对于 Killing 场运动，但不必朝着黑洞的方向运动。由式 (6.12) 和式 (6.14)，当我们趋近静止极限面时，稳态观者的红移发散，因此这个曲面也被称为\textbf{无限红移面}（infinite redshift surface）。正如我们在讨论 Kerr 度规时将会看到的，静止极限面与事件视界之间的区域——能层（ergosphere）——是一个类时路径不可避免地被拖着随黑洞一起转动的地方。在稳态时空中我们将继续使用“表面引力”这个标签，并用 Killing 矢量 $\chi^{\mu}$（它在事件视界上确实变为类光）来计算它，即使所得的量已不能解释为无穷远处所看到的静态观者的引力加速度。

让我们把这些概念应用到 Schwarzschild 情形，看看它们如何运作。对于度规
\begin{equation*}
\mathrm{d}s^{2} = -\left(1-\frac{2GM}{r}\right)\mathrm{d}t^{2} + \left(1-\frac{2GM}{r}\right)^{-1}\mathrm{d}r^{2} + r^{2}\,\mathrm{d}\Omega^{2},
\tag{6.18}
\end{equation*}

% ---- 书页 248 / p263 ----
Killing 矢量和静态四速度为
\begin{equation*}
K^\mu = (1,\ 0,\ 0,\ 0), \qquad U^\mu = \left[\left(1-\frac{2GM}{r}\right)^{-1/2},\ 0,\ 0,\ 0\right],
\tag{6.19}
\end{equation*}
于是红移因子为
\begin{equation*}
V = \sqrt{1-\frac{2GM}{r}}.
\tag{6.20}
\end{equation*}
（注意这与上一章中我们对红移的计算相符。）由式 (6.15)，加速度为
\begin{equation*}
a_\mu = \frac{GM}{r^{2}\left(1-\dfrac{2GM}{r}\right)}\nabla_\mu r,
\tag{6.21}
\end{equation*}
其中当然有 $\nabla_\mu r=\delta^{r}_{\mu}$。于是加速度的大小为
\begin{equation*}
a = \frac{GM}{r^{2}\left(1-\dfrac{2GM}{r}\right)^{1/2}}.
\tag{6.22}
\end{equation*}
表面引力是 $\kappa=Va$ 在事件视界 $r=2GM$ 处的取值，而
\begin{equation*}
Va = \frac{GM}{r^{2}},
\tag{6.23}
\end{equation*}
所以 Schwarzschild 黑洞的表面引力为
\begin{equation*}
\kappa = \frac{1}{4GM}.
\tag{6.24}
\end{equation*}
表面引力随质量增大反而减小，这看上去也许令人诧异，但只要看一眼式 (6.23) 就明白是怎么回事了：在固定半径处，增大 $M$ 确实会使组合 $Va$ 变大，但质量的增大同时也使 Schwarzschild 半径增大，后一种效应占了上风。因此，大黑洞的表面引力实际上比小黑洞的还要弱；这与对潮汐力的考察相一致——黑洞越大，潮汐力也越小。

\section{质量、电荷与自旋}

我们在上面曾宣称，广义相对论最一般的稳态黑洞解由质量、电荷和自旋刻画，那么就应该考虑这些量在 GR 中如何定义。电荷最容易考虑，所以我们从它开始；更多细节见附录 E 中关于 Stokes 定理的讨论。我们将具体考察电荷，不过磁荷也可以用同样的方式来处理。

% ---- 书页 249 / p264 ----
麦克斯韦方程组把电磁场强张量 $F_{\mu\nu}$ 与电流四矢量 $J^{\mu}_{e}$ 联系起来，
\begin{equation*}
\nabla_\nu F^{\mu\nu} = J^{\mu}_{e}.
\tag{6.25}
\end{equation*}
穿过类空超曲面 $\Sigma$ 的电荷由该超曲面上坐标 $x^{i}$ 的积分给出，
\begin{align*}
Q &= -\int_{\Sigma} d^{3}x\,\sqrt{\gamma}\, n_\mu J^{\mu}_{e}\\
  &= -\int_{\Sigma} d^{3}x\,\sqrt{\gamma}\, n_\mu\nabla_\nu F^{\mu\nu},
\tag{6.26}
\end{align*}
其中 $\gamma_{ij}$ 是与 $\Sigma$ 相联系的诱导度规，$n^{\mu}$ 是单位法矢量。这个负号保证正的电荷密度加上指向未来的法矢量会给出正的总电荷。于是可以用 Stokes 定理把电荷表示成一个边界积分，
\begin{equation*}
Q = -\int_{\partial\Sigma} d^{2}x\,\sqrt{\gamma^{(2)}}\, n_\mu\sigma_\nu F^{\mu\nu},
\tag{6.27}
\end{equation*}
其中边界 $\partial\Sigma$——典型地是空间无穷远处的一个二维球面——具有度规 $\gamma^{(2)}_{ij}$ 和指向外部的法矢量 $\sigma^{\mu}$。磁荷则可以通过把 $F^{\mu\nu}$ 换成对偶张量 ${*F}^{\mu\nu}=\frac{1}{2}\epsilon^{\mu\nu\rho\sigma}F_{\rho\sigma}$ 来确定。因此，要计算总电荷，我们只需要知道电磁场在空间无穷远处的行为。在附录 E 中，我们对 Minkowski 空间中的一个点电荷作了显式计算，得到的结果不出所料，但很好地检验了我们的约定确实自洽。

现在我们转向渐近平直时空的总能量（或质量）的概念。这是一个比电荷棘手得多的概念；一方面，在广义相对论中能量-动量是张量而不是矢量；另一方面，能量-动量张量 $T_{\mu\nu}$ 只描述物质的性质，而不描述引力场的性质。但回忆一下，第 3 章中我们讨论过：如果时空是稳态的、具有类时 Killing 矢量场 $K^{\mu}$，我们仍然可以定义一个守恒的总能量。我们先构造一个流
\begin{equation*}
J^{\mu}_{T} = K_\nu T^{\mu\nu},
\tag{6.28}
\end{equation*}
其中 $T^{\mu\nu}$ 是能量-动量张量。因为这个流是无散的（由 Killing 方程和 $T^{\mu\nu}$ 的守恒而来），我们可以在一个类空曲面 $\Sigma$ 上积分，得到一个守恒的能量，
\begin{equation*}
E_{T} = \int_{\Sigma} d^{3}x\,\sqrt{\gamma}\, n_\mu J^{\mu}_{T},
\tag{6.29}
\end{equation*}
就像电荷那样。这个表达式虽然有趣，但显然有几点不足。例如考虑 Schwarzschild 度规，它有一个
% ---- 书页 250 / p265 ----
Killing 矢量，但 $T^{\mu\nu}$ 处处为零。那么 Schwarzschild 黑洞的能量就是零吗？无论从物理还是数学的理由看，都有理由怀疑并非如此；毕竟那里有一个奇点，它使得这个积分很难算出来。此外，Schwarzschild 黑洞可以由一颗具有确定非零能量的大质量恒星演化而来，我们或许愿意让这个能量守恒。因此值得去寻找另一种能量定义，以便更好地捕捉我们对黑洞时空的直观图像。

我们暂时仍然限于具有类时 Killing 矢量 $K^{\mu}$ 的时空，考虑一个新的流
\begin{equation*}
J^{\mu}_{R} = K_\nu R^{\mu\nu}.
\tag{6.30}
\end{equation*}
利用 Einstein 方程，我们也可以把它等价地写成
\begin{equation*}
J^{\mu}_{R} = 8\pi G K_\nu\left(T^{\mu\nu}-\frac{1}{2}T g^{\mu\nu}\right).
\tag{6.31}
\end{equation*}
Ricci 张量并不是无散的；我们有的是缩并的 Bianchi 恒等式，
\begin{equation*}
\nabla_\mu R^{\mu\nu} = \frac{1}{2}\nabla^{\nu}R.
\tag{6.32}
\end{equation*}
但这个恒等式加上 Killing 方程，就足以保证我们的新流是守恒的。为看清这一点，我们只需计算
\begin{equation*}
\nabla_\mu J^{\mu}_{R} = (\nabla_\mu K_\nu)R^{\mu\nu} + K_\nu(\nabla_\mu R^{\mu\nu}).
\tag{6.33}
\end{equation*}
第一项为零，因为 $R^{\mu\nu}$ 是对称的，而 $\nabla_\mu K_\nu$ 是反对称的（由 Killing 方程）。于是利用式 (6.32)，我们得到
\begin{equation*}
\nabla_\mu J^{\mu}_{R} = \frac{1}{2}K_\nu\nabla^{\nu}R = 0,
\tag{6.34}
\end{equation*}
我们知道它为零，因为 $R$ 沿 Killing 矢量的方向导数为零，见式 (3.178)。

和前面一样，我们可以定义一个与这个流相联系的守恒能量，
\begin{equation*}
E_{R} = \frac{1}{4\pi G}\int_{\Sigma} d^{3}x\,\sqrt{\gamma}\, n_\mu J^{\mu}_{R},
\tag{6.35}
\end{equation*}
其中归一化的选取是为了将来的方便。能量 $E_{R}$ 将不依赖于类空超曲面 $\Sigma$ 的选取，因而是守恒的。这个能量概念比起 $E_{T}$ 有一项显著的优势，源于这样一个事实：$E_{R}$ 可以改写成空间无穷远处一个二维球面上的面积分。为看清这一点，回忆式 (3.177)：任何 Killing 矢量都满足 $\nabla_\mu\nabla_\nu K^{\mu}=K^{\mu}R_{\mu\nu}$；于是这个流本身可以写成一个全导数，
\begin{equation*}
J^{\mu}_{R} = \nabla_\nu(\nabla^{\mu}K^{\nu}),
\tag{6.36}
\end{equation*}
